\documentclass[11pt]{article}
\usepackage[a4paper,margin=2.6cm]{geometry}
\usepackage{amsmath,amssymb,amsthm}
\usepackage{booktabs}
\usepackage[hidelinks]{hyperref}
\usepackage{xcolor}
\usepackage{needspace}

\newtheorem{proposition}{Proposition}
\newtheorem{corollary}{Corollary}
\theoremstyle{remark}
\newtheorem{remark}{Remark}

\newcommand{\F}{\mathbb{F}}
\newcommand{\Sed}{\mathbb{S}}
\DeclareMathOperator{\rank}{rank}
\DeclareMathOperator{\col}{col}
\DeclareMathOperator{\row}{row}

% Copy-paste fidelity: map glyphs to Unicode, and in typewriter text use cmtt's real underscore glyph (slot 95)
% instead of the rule \textunderscore draws, so file names and commands copy and search correctly. A mild
% penalty lets a long identifier break after an underscore when nothing better is available.
\input{glyphtounicode}
\pdfgentounicode=1
\makeatletter
\DeclareRobustCommand{\_}{%
  \edef\@tempa{\ttdefault}%
  \ifx\f@family\@tempa\char95 \else\textunderscore\fi\penalty300\relax}
\makeatother
% DOIs: clickable, never hyphenated
\newcommand{\doi}[1]{\href{https://doi.org/#1}{\nolinkurl{#1}}}
% List bullets from the math symbol font (Type 1), not the text-companion font, which pdflatex would embed as a bitmap
\renewcommand{\labelitemi}{$\bullet$}

\title{Generator-induced invariants in lattice public keys:\\ a pre-standardization rank diagnostic}
\author{Matthew Gifford\\ \small Independent researcher, Hollister, California, USA}
\date{\small October 2026}

\begin{document}
\maketitle

\begin{abstract}
A public-key construction that builds its public object by applying a fixed family of linear maps to random
seeds---to compress the key, to inherit a symmetry, or to support a reduction---produces an object that is the image of
a single fixed \emph{design matrix}. Whatever the seeds, the rows of the public object then lie in the row space of
that matrix, so its rank can never exceed the matrix's. When the rank $\rho$ of the design matrix is below the ambient dimension, every public key is
rank-deficient: exact linear algebra distinguishes it from the uniform matrix a security reduction assumes, and
recovering the noise reduces to an LWE instance of dimension at most $\rho$. The argument takes one line---the rank of a
product is at most the rank of either factor---yet the weakness is easy to miss: a rank check on a realized instance of
dimension $n$ can certify only $\rho\ge n$, so a check run below the ceiling returns a clean result on a broken scheme.
We make this concrete with a reproducible case study, a withdrawn sedenion Module-LWE proposal whose public matrix has
rank at most 76 at every module rank and whose own structural check, run at module rank 4, could not see the ceiling.
We package the lesson as an exact and cheap diagnostic---count, then compute the rank of the design matrix over the
working field at the deployed size, and confirm it on an instance larger than the ceiling---with a dependency-free
implementation that reproduces the case study, checks the design against the scheme's own realizations, and passes
structured controls, one of them shaped like ML-KEM, and we place the mechanism among published breaks in which structure introduced for a benefit became the attack surface. Passing the diagnostic is
necessary for any security argument that assumes a uniform, full-rank matrix, but not sufficient, and it costs less than
the lattice estimate it should precede.
\end{abstract}

\paragraph{Epistemic conventions.} As in the author's other work, claims are tiered. Every proposition and corollary
below is T1: proven from the stated definitions, or verified by exact computation. The remarks are held to the same
standard: they state only what follows from the propositions, from the cited literature, or from exact computation that
the accompanying self-test checks. Section~\ref{sec:pattern} classifies published breaks; that classification is
interpretive, and where it reaches beyond Proposition~\ref{prop:ceiling} it says so.

\section{Introduction}\label{sec:intro}

Uniformly random public keys are large, and designers routinely replace them with a \emph{generated} object that is
cheaper to store or transmit: the cyclic and negacyclic structure of Ring- and Module-LWE, the quasi-cyclic and dyadic
codes of compact code-based schemes, group-action images, and orbit constructions of many kinds. The generator is always
introduced for a reason: a smaller key, a symmetry the designer believes is benign or even helpful, or a structural
argument that the construction reduces to a well-studied problem.

The thesis of this note is that such a generator is a fixed linear gadget, and that random seeds fed through a fixed
linear gadget cannot produce anything outside its image. If that image is confined---a rank ceiling, equivalently a
kernel common to every key---the public object inherits the confinement however much entropy the seeds carry. When the
ceiling is below the ambient dimension, the consequences are immediate: the public key is distinguishable from the
uniform matrix a reduction assumes, and recovery drops from the ambient dimension to at most the rank of the design
matrix. Not every invariant of a generator is a weakness: the multiplication maps of Ring- and Module-LWE share fixed
subspaces, and no attack on ML-KEM exploits them; Remark~\ref{rem:inv} says what separates the two.

None of the linear algebra is deep, and that is the point. The mechanism is elementary, and the diagnostic that catches
it takes seconds of exact arithmetic. Yet the weakness is easy to ship, because a rank check on a realized instance
certifies only that the ceiling is at least the size of that instance. Below the ceiling, a generated object can be full
rank and, by its rank, indistinguishable from a uniform control; a check run there---the diligent thing to do---returns a clean
result. A scheme can therefore pass its own structural test and still carry a fatal ceiling, if the test was run on too
small an instance.

\paragraph{Contributions.}
\begin{enumerate}
\item A general statement (Section~\ref{sec:mech}): a seed-generated public object factors as $M=R\cdot D$ with $D$ a
  fixed design matrix, so $\rank M\le\rank D$ for every instance size and every draw of the seeds. The same
  factorization gives a kernel shared by every public key (a distinguisher) and confines the LWE error to an affine
  subspace of dimension at most $\rank D$ (a reduced recovery problem), and two counts bound $\rank D$ before anything
  is computed.
\item A visibility principle (Proposition~\ref{prop:vis}): a full-rank check at ambient dimension $n$ certifies only
  $\rank D\ge n$, and a deficiency is guaranteed to show only once $n>\rank D$. Computing $\rank D$ from the design
  certifies every size at once.
\item A reproducible case study (Section~\ref{sec:case}): a withdrawn sedenion Module-LWE proposal in which exactly this
  mechanism gives a ceiling of 76 at every module rank, with an account of why the scheme's own check, run at module
  rank 4, could not see it.
\item An exact, cheap diagnostic (Section~\ref{sec:diag}) with a dependency-free implementation that reproduces the case
  study as a regression test, takes its verdict at the deployed size, checks the design against the scheme's
  realizations, and passes structured controls: a sedenion module matrix with independent uniform entries, an
  ML-KEM-shaped module lattice, and a row-seeded design with a long enough seed.
\item A short classification of published breaks of the same shape (Section~\ref{sec:pattern}), with an account
  (Section~\ref{sec:limits}) of what the theorem does and does not cover.
\end{enumerate}

\paragraph{Scope.} This is a defensive, methodological note. The case study is the author's own withdrawn proposal, and
the ML-KEM design~\cite{FIPS} appears only as a correctly designed control, which the diagnostic passes. No attack on
any deployed or standardized system is given or implied.

\section{The mechanism}\label{sec:mech}

\subsection{Setup}\label{sec:setup}

Let $\F$ be a field. A \emph{seed-generated public object} is built as follows. The scheme fixes a finite family of
$\F$-linear maps $T_0,\dots,T_{r-1}\colon\F^d\to\F^d$ (the \emph{generator}), draws independent random seeds, applies
the generator to the seeds, and combines the results $\F$-linearly. A \emph{seed} here is the vector of coefficients the
generator consumes, after any nonlinear expansion: ML-KEM derives its matrix from one 32-byte string, but its generator
consumes $256k^2$ coefficients~\cite{FIPS}. Write $M$ for the matrix of the $\F$-linear map the scheme exposes---for an
LWE-type key, the map $s\mapsto As$. Because the assembly is linear in the seeds and linear in the fixed generator
images, $M$ is bilinear in the two and factors as
\[
  M \;=\; R\cdot D,
\]
where $R\in\F^{a\times g}$ carries all of the seed randomness (its entries are linear in the seed coordinates) and
$D\in\F^{g\times b}$ is \emph{fixed}: its entries are determined by the generator and the combination rule and do not
depend on the seeds. The factorization is not unique: $R=M$ with $D$ the identity always works, so a full-rank $D$
proves nothing by itself. Write $M(\theta)$ for the public map at the seed vector $\theta$, which is linear in $\theta$,
and let $W=\sum_m\row M(e_m)$, summed over the unit seed vectors $e_m$. Every row of every $M(\theta)$ lies in $W$, and
every valid $D$ has $\row(D)\supseteq W$. We call $D$ the \emph{design matrix} when $\row(D)=W$, the minimal choice.
Only its row space matters below, so repeated rows may be dropped. A larger valid $D$ still bounds the rank, so a
ceiling it shows is real, but it can hide one; a pass means something only for the minimal design. For a uniform object
$W$ is the whole space and $D$ is the identity.

\subsection{The rank ceiling}

\begin{proposition}[rank ceiling]\label{prop:ceiling}
For every instance and every draw of the seeds, $\rank M\le\rank D$.
\end{proposition}

\begin{proof}
$M=RD$ and $\row(RD)\subseteq\row(D)$, so $\rank M\le\rank D$.
\end{proof}

The content lies entirely in the fact that $D$ is fixed. Seed entropy moves $M$ within the row space of $D$; it cannot
raise the rank above $\rank D$. If $\rank D$ is below the ambient dimension, every public key the scheme can emit is
rank-deficient, whatever the seeds.

Two counts bound $\rank D$ before anything is computed: the number of distinct rows of $D$ and the number of distinct
columns. When each block row of $M$ reuses one seed of length $\ell$ through the generator, as in the case study, the
design needs only $\ell d$ rows for $d\times d$ blocks, whatever the generator; and for a cyclic generator of order $r$
it has at most $rd$ distinct columns (Section~\ref{sec:vis}). If either count is below the number of columns of $M$ at
the deployed size, the scheme fails without further computation.

\begin{corollary}[distinguisher]\label{cor:dist}
Let $\rho=\rank D<b$. Then $\ker D\subseteq\ker M$ for every public key, and $\dim\ker D=b-\rho$. In particular, a fixed
nonzero $x\in\ker D$---computable from the design alone---satisfies $Mx=0$ for every public key. A uniform
$M\in\F_q^{a\times b}$ does so with probability $q^{-a}$, and so does a module matrix whose entries are independent
uniform elements of a $d$-dimensional algebra, provided right multiplication by some block $x_j$ of $x$ is invertible.
\end{corollary}

\begin{proof}
If $Dx=0$ then $Mx=R(Dx)=0$; the dimension count is rank--nullity. For uniform $M$, each of the $a$ rows is orthogonal
to $x\neq0$ independently with probability $1/q$. For independent uniform entries $A_{ij}$, block row $i$ of $Mx$ is
$\sum_j A_{ij}x_j$; the term $A_{ij}x_j$ is uniform when right multiplication by $x_j$ is invertible, so the block row
vanishes with probability $q^{-d}$, independently over the $a/d$ block rows.
\end{proof}

A structured matrix is distinguishable from a uniform $M\in\F_q^{a\times b}$ whether or not it is weak---ML-KEM's is,
through its negacyclic blocks---so the comparison that matters is with the distribution the scheme's security argument
assumes. The corollary covers both cases; Section~\ref{sec:ceiling} applies it to the case study.

\begin{corollary}[recovery]\label{cor:rec}
Write the public key in LWE form $t=Ms+e$ with $M\in\F_q^{a\times b}$, and let $V=\col(M)$ and $\rho_M=\dim V\le\rank D$.
Then $t-e\in V$, so $e$ lies in the known affine subspace $t+V$ of dimension $\rho_M$. Choosing $\rho_M$ coordinates $I$
on which a basis of $V$ is invertible expresses the remaining $a-\rho_M$ coordinates of $e$ as fixed affine functions of
$e_I$: an LWE instance with secret $e_I$ of dimension $\rho_M$ and $a-\rho_M$ samples. Equivalently, a parity-check
matrix $L$ for $V$, with $LM=0$, gives $Lt=Le$.
\end{corollary}

\begin{proof}
Every $v\in V$ is determined by $v_I$, with $v_{\bar I}=Cv_I$ for a matrix $C$ computable from $M$. Applied to
$v=t-e$ this gives $e_{\bar I}=Ce_I+(t_{\bar I}-Ct_I)$, in which $e_I$ and $e_{\bar I}$ are short.
\end{proof}

\paragraph{From noise to message.} Recovering the noise suffices to decrypt whenever decryption is linear. Suppose a
ciphertext is $u=M^{*}r+e_1$ and $v=\langle t,r\rangle+e_2+\mathrm{encode}(m)$, where $M^{*}$ is the adjoint of $M$ for
the pairing used in decryption and $r$, $e_1$, $e_2$ are short. Since $\rank M^{*}=\rank M$, the argument of
Corollary~\ref{cor:rec} applied to $u$ reduces finding $e_1$ to an LWE instance of the same dimension as for $e$. Once
both are found, any $s'$ with $Ms'=t-e$ gives
$\langle s',u-e_1\rangle=\langle Ms',r\rangle=\langle t-e,r\rangle$, so
$v-\langle s',u-e_1\rangle=\langle e,r\rangle+e_2+\mathrm{encode}(m)$, a small perturbation of the encoded message.

\begin{remark}\label{rem:inv}
Rank is the cheapest invariant to read off, and the only one this argument transfers: the factorization gives
$\row(M)\subseteq\row(D)$, equivalently $\ker D\subseteq\ker M$, and nothing more. A subspace that every key preserves is
a different matter, because it need not be a weakness. Over $\F_{3329}$ the polynomial $X^{256}+1$ splits into 128
quadratic factors, so every multiplication map of ML-KEM's ring preserves the same 128 two-dimensional subspaces (its
NTT slots); over $\F_{8380417}$ it splits into linear factors, so the multiplication maps of ML-DSA share 256
eigenvectors~\cite{FIPS}. Projecting onto such a subspace does shrink the secret, yet no attack on either scheme
exploits it, because the projected noise is close to uniform. An invariant subspace is a weakness only when the projected
noise is far from uniform: short, as under the restriction to coordinates in Corollary~\ref{cor:rec}, or confined to a
small set, as under evaluation at a suitable root of the defining polynomial (for instance $1$) in the provably weak
instances of Ring-LWE~\cite{ELOS15,CIV16}; Peikert~\cite{Pei16} shows that those attacks do not apply when the error is
wide enough relative to the ring's geometry. A test for invariant subspaces must check that condition, not the subspace
alone, and this note's diagnostic does not attempt it.
\end{remark}

\subsection{Why the ceiling hides}\label{sec:vis}

The design matrix of an instance is assembled from the generator images that instance uses. Suppose the generator
contributes $c$ independent images in total; for a cyclic generator of order $r$ acting on a $d$-dimensional alphabet,
$c\le rd$, and the column space of the design stops growing once an instance uses all $r$ powers. Let $D$ be the
\emph{full design}, which contains every generator image, and let $\rho=\rank D$. Let $M_k=R_kD_k$ be the realized public
map at instance size $k$, and assume that $D_k$ is assembled from columns of $D$, so that $\col(D_k)\subseteq\col(D)$ and
$\rank D_k\le\rho$. This holds when each block row is generated from a seed of fixed length, so that the rows of $D_k$
do not depend on $k$, and the generator has finitely many images, as in the case study; it fails for independently
seeded blocks, whose design grows with $k$. Let $n(k)$ be the smaller dimension of
$M_k$, which is the rank a uniform matrix of that shape has with high probability.

\begin{proposition}[visibility]\label{prop:vis}
\leavevmode
\begin{enumerate}
\item[(i)] If $\rank M_k=n(k)$, then $\rho\ge n(k)$, and Proposition~\ref{prop:ceiling} excludes no larger value of
  $\rho$: a full-rank result at size $k$ is consistent with every ceiling $\rho\ge n(k)$, including ceilings far below
  the ambient dimension $n(k')$ of a larger deployed size $k'$.
\item[(ii)] If $n(k)>\rho$, then $\rank M_k\le\rho<n(k)$: the deficiency is visible.
\end{enumerate}
\end{proposition}

\begin{proof}
Proposition~\ref{prop:ceiling} applied to $M_k=R_kD_k$, together with $\rank D_k\le\rho$ and the shape of $M_k$, gives
$\rank M_k\le\min(n(k),\rho)$. If $\rank M_k=n(k)$, this forces $n(k)\le\rho$ and constrains $\rho$ no further.
Table~\ref{tab:rank} exhibits a design with $\rho=76$ whose realizations have full rank at every $k\le4$ ($n(k)\le64$),
although the deployed size has $n(k')=1024$. Part~(ii) is immediate.
\end{proof}

\begin{remark}\label{rem:tight}
In the case study of Section~\ref{sec:case}, the realized rank equals $\rank N_k$, the rank of the design restricted to
the powers $\sigma^0,\dots,\sigma^{k-1}$ that an instance of size $k$ uses, at every size tested (Table~\ref{tab:rank}):
$16k$ for $k\le4$, then 72 and 74 at $k=5$ and $6$, and $\rho=76$ from $k=7$, once all seven powers appear
($\sigma^7=\mathrm{id}$). Proposition~\ref{prop:ceiling} is therefore tight here for generic seeds. The deficiency first
appears at $k=5$, the smallest $k$ with $n(k)>\rho$, and a check at $k=4$ is precisely the uninformative case~(i). This
is a property of the case-study design, established by exact computation and checked by the self-test; it is not a
theorem about all designs.
\end{remark}

Two practical rules follow. An empirical rank check certifies only the size it is run at, so a check below the deployed
size certifies only a lower bound on the ceiling, not the deployed key. And computing $\rho=\rank D$ from the design certifies every size at once:
it is cheaper and stronger than any empirical check, which then serves only to confirm that the factorization was done
correctly.

\section{Case study: a sedenion Module-LWE proposal}\label{sec:case}

We instantiate the mechanism on a concrete, withdrawn scheme. The numerical claims in this section are those of the
negative-result note that broke it~\cite{GifNote}, to which we refer for the computations; the aim here is to exhibit the
general mechanism, not to re-derive the break. The scheme and the break are the author's own.

\subsection{The construction and its intent}

The scheme, SLWE~\cite{GifFLT}, is a Module-LWE key encapsulation mechanism~\cite{Regev,LS15} over the 16-dimensional
sedenion algebra $\Sed_q=\Sed\otimes\F_q$, the fourth Cayley--Dickson algebra. Its public matrix
$A\in\Sed_q^{k\times k}$ is built from orbits of an order-7 coordinate permutation $\sigma$---a ``Singer cycle''
inherited from the symmetry of the Fano plane---as $A_{ij}=\sigma^{j}(\mathrm{seed}_i)$, with independent uniform seeds
$\mathrm{seed}_i\in\Sed_q$. The orbit structure was deliberate: it compresses the key (each row comes from one seed) and
was meant to preserve the algebra's zero-divisor interface. Non-associativity was argued to remove an ideal-lattice
attack class and so to reduce security to that of a generic module lattice. Two parameter sets were published with a
claim of NIST Category~5: $n=512$ ($k=32$) with $q\approx2^{32}$, and $k=64$ ($n=1024$) with $q=911$.

\subsection{The factorization and the ceiling}\label{sec:ceiling}

Let $M\in\F_q^{16k\times16k}$ be the matrix of $s\mapsto As$, so that $M_{16i+u,\,16j+c}$ is coefficient $u$ of
$\sigma^{j}(\mathrm{seed}_i)\,e_c$. Concretely, $\sigma$ fixes $e_0$ and $e_8$, maps $e_t\mapsto e_{t+1}$ for
$1\le t\le6$ and $e_7\mapsto e_1$, and acts in the same way on $e_9,\dots,e_{15}$; the product is the Cayley--Dickson
doubling $(a,b)(c,d)=(ac-\bar d b,\;da+b\bar c)$. Because $\sigma$ is linear and the sedenion product is bilinear, $M$
factors exactly as in Section~\ref{sec:mech}:
\[
  M=(S\otimes I_{16})\,N_k,
\]
where $S\in\F_q^{k\times16}$ holds the seeds (this is $R$) and $N_k$ is assembled from the columns of a fixed integer
matrix $N\in\mathbb{Z}^{256\times112}$ whose entry $N_{16m+u,\,16j+c}$ is coefficient $u$ of $\sigma^{j}(e_m)\,e_c$
(this is $D$). Since $\sigma^7=\mathrm{id}$, the column blocks of $N_k$ repeat with period 7, so $\col(N_k)\subseteq\col(N)$,
with equality once $k\ge7$.

Counting alone already bounds the rank. The design $N$ has 256 rows---16 seed coordinates times 16 rows per block---and
112 columns---seven powers of $\sigma$ times 16---so $\rank M\le112$ at every $k$, and for $k>7$ the repetition
$A_{i,j+7}=A_{ij}$ is visible in the deployed matrix by inspection. Exact computation sharpens 112 to 76: $\rank N=76$
over $\mathbb{Q}$ and modulo $911$ and $q=4\,294\,977\,961$~\cite{GifNote}, the least prime $\equiv1\pmod{455}$ above
$2^{32}$, with which~\cite{GifNote} instantiates the first set's $q\approx2^{32}$. Because $N$ has integer entries, its rank
modulo any prime is at most its rank over $\mathbb{Q}$, so the ceiling of 76 holds at every modulus at once. Over
$\F_{911}$, which contains the seventh roots of unity, a discrete Fourier transform over the power index $j$ splits
$\col(N)$ into a part of dimension 16 for the trivial character and six parts of dimension 10, so $76=16+6\cdot10$.
The other common Cayley--Dickson convention, $(a,b)(c,d)=(ac-d\bar b,\;\bar a d+cb)$, defines the opposite algebra, in
which $xy$ becomes $yx$, and gives the same ranks. The accompanying diagnostic re-derives $\rank N=76$ over $\mathbb{Q}$ by fraction-free elimination
and modulo $911$, $2731$, $2^{31}-1$ and $q$. By Proposition~\ref{prop:ceiling}, $\rank M\le76$ at every module rank,
against $16k$ for a matrix with uniform entries. Table~\ref{tab:rank} lists the exact ranks.

\begin{table}[h]\centering\small
\begin{tabular}{lccccccccccc}\toprule
$k$ & 1&2&3&4&5&6&7&8&16&32&64\\\midrule
$n(k)=16k$ & 16&32&48&64&80&96&112&128&256&512&1024\\
$\rank M_k$ (Singer orbit) & 16&32&48&64&\textbf{72}&\textbf{74}&\textbf{76}&\textbf{76}&\textbf{76}&\textbf{76}&\textbf{76}\\
$\rank M_k$ (uniform entries) & 16&32&48&64&80&96&112&128&256&512&1024\\\bottomrule
\end{tabular}
\caption{Exact rank of the flattened public map at module rank $k$: Singer orbit versus a matrix with independent
uniform sedenion entries, the distribution the scheme's reduction assumed (a uniform matrix over $\F_q$ has the same
ranks). The values over $\F_{911}$ and over $\F_q$ with $q=4\,294\,977\,961$ are identical; the Singer-orbit row is as
reported in~\cite{GifNote}. The self-test of the accompanying diagnostic (Section~\ref{sec:diag}) regenerates the columns
$k\le8$ over $\F_{911}$ and checks that each Singer-orbit entry equals $\rank N_k$. Bold entries fall short of $n(k)$.
The columns $k\le4$ are identical for the two matrices: a check there cannot see the ceiling
(Proposition~\ref{prop:vis}(i)).}\label{tab:rank}
\end{table}

By Corollary~\ref{cor:rec}, recovering the noise is an LWE instance of dimension 76 with $16k-76$ samples. The lattice
estimator~\cite{APS15,estimator} places it at $2^{38.6}$ operations at the first parameter set (436 samples) and at
$2^{40.4}$ at the second (948 samples), with the block size at the floor of its model, so that reduction at the scale of
LLL suffices~\cite{GifNote}. The encryption map is the conjugate transpose of $A$~\cite{GifNote}, of the same rank, so
each ciphertext's noise reduces in the same way, and the argument after Corollary~\ref{cor:rec} then yields the message;
\cite{GifNote} derives this step. Distinguishing the public
matrix from the distribution its reduction assumed is cheaper still: by Corollary~\ref{cor:dist}, it takes Gaussian
elimination. A fixed $x\in\ker N_8$ is annihilated by every Singer-orbit key at $k=8$, while right multiplication by one
of its blocks is invertible, so a matrix with independent uniform sedenion entries annihilates it only with probability
$q^{-128}$; the self-test checks both facts over $\F_{911}$.

\subsection{Why the scheme's own check missed it}\label{sec:missed}

The proposal included a structural self-test: a five-test Gram--Schmidt profile of the lattice, looking for exactly the
kind of anomaly at issue. It reported no rank deficiency, no anomalous short vectors and no invariant sublattice
relative to a uniform matrix~\cite{GifFLT}. It was run on a $64\times64$ proxy at module rank $k=4$, and the proposal
itself labelled it a toy-scale result that ``does not establish anything at production dimension''.

Proposition~\ref{prop:vis} makes that caveat exact. At $k=4$ the public map has $n(4)=64$ columns, below the ceiling of
76, so by part~(i) a full-rank result certified only $\rho\ge64$. An instance of size 4 also uses only four of the seven
powers of $\sigma$, so even the repetition $A_{i,j+7}=A_{ij}$ of the deployed matrix was absent from it. The deployed
matrix at $k=64$ has $n=1024$; every value $64\le\rho<1024$ was left open, and the truth was $\rho=76$. The check returned
a clean result not because the scheme was sound but because it was run below the ceiling. Table~\ref{tab:rank} shows the
trap at a glance: the columns $k\le4$ are identical for the Singer-orbit and the uniform matrix, and they separate only
from $k=5$.

\subsection{Non-associativity was beside the point}

The scheme's security narrative rested on the non-associativity of the sedenions and on the $\sigma$-symmetry of $A$.
Neither bears on the break. Decryption uses the bilinear conjugate-norm pairing, so every map the attacker needs is
$\F_q$-linear, and the argument of Section~\ref{sec:mech} uses only linearity. The permutation $\sigma$ is not an
automorphism of the sedenion product~\cite{GifNote} but a relabelling of coordinates: it breaks
$\sigma(xy)=\sigma(x)\sigma(y)$ on 144 of the 256 products of basis elements. It is not even a collineation of the Fano plane as the product labels
it, whose lines are $\{1,2,3\}$, $\{1,4,5\}$, $\{1,6,7\}$, $\{2,4,6\}$, $\{2,5,7\}$, $\{3,4,7\}$ and $\{3,5,6\}$: it
sends $\{1,2,3\}$ to $\{2,3,4\}$. So the ``symmetry'' of $A$ is not an algebraic symmetry at all. The lesson travels:
advertised algebraic richness can be orthogonal to the property that carries security, and a security argument phrased
in terms of the rich structure may be answering the wrong question.

\subsection{The repair that was not one}\label{sec:repair}

Once the deficiency was found, the natural repair was to replace the Singer-orbit matrix with a uniform one, which
removes the ceiling. It does not rescue the claim. At the first parameter set the modulus-to-noise ratio alone puts the
cheapest attack near $2^{51}$. At the second, hybrid attacks that exploit the weight-64 sparse secret put it near
$2^{167}$, about 95 bits below the $2^{262}$ the same estimator gives ML-KEM-1024, although primal uSVP alone costs
about $2^{254}$~\cite{GifNote}. Clearing the structural flag one has found does not clear a parameter problem one has
not. Passing the rank test of Section~\ref{sec:diag} is necessary for a uniform-matrix security argument, never
sufficient.

\section{The diagnostic}\label{sec:diag}

The mechanism gives a checklist that a designer or evaluator can run before trusting a compressed or structured public
key, at a cost far below the lattice estimate it should precede.
\begin{enumerate}\setcounter{enumi}{-1}
\item \textbf{Count.} Count the distinct rows and the distinct columns of the design at the deployed size; $\rank D$ is
  at most the smaller. If that is below the number of columns of $M$, the scheme fails before any elimination (in the
  case study, 112 against 512 and 1024).
\item \textbf{Factor.} Write the construction as $M=R\cdot D$, separating the seed randomness $R$ from the fixed
  generator gadget $D$, and take $D$ minimal: its row space must be spanned by the realizations at unit seeds
  (Section~\ref{sec:setup}). For a uniform object $D$ is the identity, and the test passes.
\item \textbf{Build the design} at each deployed size and, when the hypothesis of Section~\ref{sec:vis} holds, the full
  design that contains every generator image (for a cyclic generator of order $r$, all $r$ powers), which bounds every
  size at once.
\item \textbf{Compute the ceiling exactly.} Compute $\rho=\rank D$ over the working field $\F_q$ with exact arithmetic,
  not floating point. When $D$ has integer entries, also compute its rank over $\mathbb{Q}$: the rank modulo any prime
  is at most that, so it bounds every modulus at once (Proposition~\ref{prop:ceiling}).
\item \textbf{Confirm on a large enough instance.} Realize $M$ at the smallest size at which $\rank D_k$ falls below the
  number of columns, and compare it, over several seed sets, with matrices of the same shape drawn from the distribution
  the scheme's security argument assumes. A check below that size is uninformative (Proposition~\ref{prop:vis}); a
  realized row outside $\row(D_k)$ contradicts Proposition~\ref{prop:ceiling} and means the factorization is wrong.
\item \textbf{Flag.} At each deployed size, report a distinguisher if $\rho$ is below $b$, the number of columns of $M$
  (Corollary~\ref{cor:dist}), and a ceiling if $\rho$ is also below the smaller dimension of $M$, with a recovery problem
  of dimension at most $\rho$ (Corollary~\ref{cor:rec}) at which the lattice estimate must be rerun. For a square map the
  two coincide; a wide map can have a fixed kernel while every key has full row rank. Invariant subspaces are not
  flagged: they are a weakness only under the condition of Remark~\ref{rem:inv}, which this test does not check.
\end{enumerate}
The test is exact linear algebra on a matrix the size of the design: it takes seconds, needs no lattice reduction or
estimator, and costs less than the Core-SVP estimate a candidate already runs. It belongs \emph{before} that estimate,
because a positive result changes the dimension at which the estimate must be run and, when $\rho$ is as small as in the
case study, makes it moot.

\paragraph{Where to look.} The factorization of step~1 exists whenever the public object is a generator applied to
seeds: cyclic or orbit-generated matrices, images under linear group actions, structured-code substitutions, and any
object described as compressed, symmetric, or inheriting an algebraic structure.

\paragraph{Implementation.} A standalone implementation in pure Python, with no dependencies, accompanies this note
(\texttt{design\_matrix\_rank.py}; Section~\ref{sec:repro}). A user describes a construction by implementing four
methods---\texttt{design\_matrix(k)}, \texttt{realize(k, seeds, p)}, \texttt{ambient\_rank(k)} and
\texttt{random\_seeds(k, p, rng)}---and listing its parameter sets, each a deployed size with its working prime; optional
methods give the shape of a non-square map, the full design, and a sampler for the distribution the security argument
assumes. On the designs the user supplies (steps 1 and 2) the tool runs steps 0, 3, 4 and 5 and takes its verdict at
each deployed size. It also checks the design against the scheme's realizations, and reports an error if a realized row
falls outside $\row(D_k)$, if the realizations are not linear in the seeds, or if, at the sizes it can afford, the
realizations at unit seeds span less than $\row(D_k)$, which would make a pass meaningless. It ships with the case study
of Section~\ref{sec:case} (verdict \emph{fail}: ceiling 76 at both parameter sets, 112 by counting alone, and the design
minimal at every size checked) and with two controls (verdict \emph{pass}): a sedenion module matrix with independent
uniform entries and an ML-KEM-shaped module of independent uniform elements of $\mathbb{Z}_{3329}[X]/(X^{256}+1)$ at
$k=2,3,4$, whose designs are full rank by construction. A \texttt{--self-test} mode reproduces the columns $k\le8$ of
Table~\ref{tab:rank}, $\rank N=76$ and the other facts this note computes; checks the verdicts on the case study, on the
controls, on a row-seeded design whose seed is long enough to pass, and on inputs built to fail or to be rejected; and
exits with an error on any mismatch, so it doubles as a regression test pinned to the published result. The empirical curve it prints marks every size below the ceiling as one at which the
collapse is invisible: Section~\ref{sec:missed} in executable form.

\section{The pattern recurs}\label{sec:pattern}

The case study is one instance of a recurring shape: \emph{a structural feature introduced for a benefit carries an
invariant, or an auxiliary datum, that an attacker reads off.} Only the first item below is an instance of
Proposition~\ref{prop:ceiling}. The others are published breaks, cited as established results; what we add is the
classification---that each turns a structure introduced for a benefit into the lever---which is interpretive, not a
theorem. All are breaks of proposed or withdrawn constructions or instantiations.
\begin{itemize}
\item \textbf{A rank ceiling from a generator.} The case study of Section~\ref{sec:case}. Benefit sought: a compressed
  key and an inherited symmetry. Lever: the rank of the design matrix. This is the linear-image case that
  Proposition~\ref{prop:ceiling} governs directly.
\item \textbf{Structured codes in McEliece.} No structural attack is known on McEliece with binary Goppa codes, but
  substituting a more structured code to shrink the key has failed repeatedly: generalized Reed--Solomon
  codes~\cite{SS92}, compact quasi-cyclic and dyadic variants~\cite{FOPT10}, and algebraic-geometry codes~\cite{CMP14},
  each broken by a structural attack on the public generator matrix that recovers the hidden code or, for
  algebraic-geometry codes, an efficient decoder. Structure added for size became the lever. Dimension counts are a
  common tool, as in Section~\ref{sec:diag}: the attack on algebraic-geometry codes works with componentwise products of
  the public code, whose dimensions are abnormally small for such codes~\cite{CMP14}. Even binary Goppa codes are
  distinguishable from random codes at high rate, where a linear system built from the public key has an abnormally
  large solution space~\cite{FGOPT13}; that result defeats an indistinguishability assumption rather than recovering a
  key.
\item \textbf{Layers in multivariate signatures.} Rainbow, a finalist in the third round of NIST's process, layered
  Unbalanced Oil and Vinegar for efficiency. The layers make some linear combinations of the public quadratic forms
  abnormally low in rank, the lever of MinRank attacks. Beullens' rectangular MinRank attack~\cite{Beu21}, and his key
  recovery for the second-round SL1 parameters in about a weekend on a laptop~\cite{Beu22}, both exploit the layered
  structure.
\item \textbf{Overstretched NTRU.} With a large modulus at fixed noise, a dense sublattice appears, which subfield and
  sublattice attacks exploit~\cite{ABD16,KF17}; the onset was later located as the ``fatigue point''~\cite{DvW21}. Here
  the structure comes from a parameter choice made for functionality rather than from an explicit generator, but the
  lever is again a fixed low-dimensional structure in the public object. It is also the published counterpart of the
  first parameter set of Section~\ref{sec:repair}, where a modulus too large for its noise, not a generator, was the
  lever.
\item \textbf{Principal ideals and short generators.} Compact key constructions over principal ideals, such as
  Soliloquy~\cite{CGS14}, admit short-generator recovery~\cite{CDPR16}: the ideal structure chosen for compactness is
  the route to the key.
\item \textbf{Weak instances of Ring-LWE.} Every multiplication map of a Ring-LWE ring preserves the same subspaces. For
  some defining polynomials and error distributions, projecting onto one of them leaves the error far from uniform and
  breaks the instance~\cite{ELOS15,CIV16}; with errors wide enough relative to the ring's geometry the projection is
  harmless~\cite{Pei16}. This is the invariant-subspace case of Remark~\ref{rem:inv}, together with the condition it
  needs.
\item \textbf{SIDH and SIKE: auxiliary data for functionality (an analogy).} The images $\phi(P),\phi(Q)$ of torsion
  points are published so that the key exchange can complete. The 2022 attacks, all published at EUROCRYPT
  2023~\cite{CD23,MMPPW23,Rob23}, use Kani's theorem~\cite{Kani97} to embed the secret isogeny in a higher-dimensional
  abelian variety and turn exactly that auxiliary data into key recovery, in polynomial time in full
  generality~\cite{Rob23}. This is not a rank ceiling and not an instance of Proposition~\ref{prop:ceiling}; it marks
  the boundary of the linear-image argument while sharing its moral: the structure published to make the scheme work
  was the attack surface.
\end{itemize}
The question this classification recommends asking of any structured proposal is simple: \emph{what did the structure I
added give the attacker?} For the linear-image cases, Section~\ref{sec:diag} answers it mechanically and cheaply; for the
others it is a design discipline rather than a single test.

\section{Limitations}\label{sec:limits}

\begin{itemize}
\item \textbf{Linear-image generators only.} Proposition~\ref{prop:ceiling} is exact and general for public objects
  bilinear in the seeds and a fixed generator. Nonlinear generators, and the SIDH case of Section~\ref{sec:pattern}, lie
  outside it; there the comparison is an analogy, and Section~\ref{sec:pattern} says so.
\item \textbf{Necessary for a uniform-matrix argument, not sufficient.} A full-rank minimal design shows only that this
  class of weakness is absent, and the tool can confirm minimality only at the sizes it realizes. Invariant subspaces,
  which a ring-structured design has whenever its defining polynomial factors modulo $q$, are a weakness only under the
  condition of Remark~\ref{rem:inv}, which the test does not check. Section~\ref{sec:repair} is the cautionary case: with
  the structure removed, the scheme still failed on its parameters---the modulus-to-noise ratio at the first set, the
  sparse secret at the second---which this test does not judge.
\item \textbf{Structure is not the enemy.} Module-LWE and binary-Goppa McEliece are structured, no structural attack on
  them is known, and the ML-KEM-shaped control passes the diagnostic. The claim is narrow and testable: generator-induced invariants must be
  computed, over the working field and for the design at the deployed size, before a structured public object is
  trusted.
\item \textbf{Dimension, not cost.} Corollary~\ref{cor:rec} gives the dimension to which recovery reduces; the concrete
  cost is a lattice-estimator question that depends, as Section~\ref{sec:repair} shows, on parameters the rank test does
  not see. A ceiling close to the ambient dimension may leave the reduced problem hard; the case study is fatal because
  $\rho=76$.
\end{itemize}

\Needspace{12\baselineskip}  % keep the heading, its lead-in and the table together
\section{Reproducibility}\label{sec:repro}

The diagnostic is a single file in the public repository \texttt{gifgaf0/gifgaf0.github.io}, and a copy accompanies
this note on Zenodo:
\begin{center}\small
\begin{tabular}{@{}ll@{}}
file & \texttt{generator\_induced\_invariants/diagnostics/design\_matrix\_rank.py}\\
SHA-256 & \texttt{24a0d4c6a40453301fc9ba7119b2d69a211dbbd67d317b0eb422e4e8ba54e876}\\
self-test & \texttt{python3 design\_matrix\_rank.py --self-test}
\end{tabular}
\end{center}
It runs under Python~3.8 or later with no other dependencies. The self-test regenerates the columns $k\le8$ of
Table~\ref{tab:rank} over $\F_{911}$ and checks that they equal $\rank N_k$; it re-derives $\rank N=76$ over $\mathbb{Q}$
by fraction-free elimination and modulo $911$, $2731$, $2^{31}-1$ and $q=4\,294\,977\,961$, and over $\mathbb{Q}$ in
the opposite algebra; it checks the count of 112, the split $76=16+6\cdot10$, the 144 failing basis products, the image
of $\{1,2,3\}$ under $\sigma$ and the kernel test of Section~\ref{sec:ceiling}. It then checks the verdicts: fail for the
case study, with its design minimal at the sizes checked; pass for both controls (the ML-KEM-shaped one at ring degree
16, with its design rank also checked at degree 256) and for a row-seeded design with a long enough seed; fail for three
designs built to fail, one of them a wide map with a fixed kernel and no rank ceiling; and an error for four inputs that
do not describe their construction, among them the case study's realization paired with the identity as its design, and
for a composite modulus. The full break of the case-study scheme---both independent rank implementations, one written blind from a
textual specification, the exact ranks computed with python-flint, and the lattice-estimator runs---is archived
with~\cite{GifNote}.

\paragraph{Acknowledgements.} AI agents (Anthropic's Claude) drafted this note, generalized its mechanism from the case
study, and wrote and ran the accompanying diagnostic; a review by Google's Gemini suggested the seed-length count of
Section~\ref{sec:mech}. The author directed the work, reviewed the text, and takes responsibility for it.

\begin{thebibliography}{99}\small
\bibitem{ABD16} M.~R.~Albrecht, S.~Bai, L.~Ducas. A subfield lattice attack on overstretched NTRU assumptions:
  cryptanalysis of some FHE and graded encoding schemes. \emph{CRYPTO} 2016. IACR ePrint 2016/127.
\bibitem{APS15} M.~R.~Albrecht, R.~Player, S.~Scott. On the concrete hardness of Learning with Errors.
  \emph{J.\ Math.\ Cryptol.} 9(3) (2015), 169--203.
\bibitem{estimator} M.~R.~Albrecht et al. The lattice estimator. \url{https://github.com/malb/lattice-estimator}.
\bibitem{Beu21} W.~Beullens. Improved cryptanalysis of UOV and Rainbow. \emph{EUROCRYPT} 2021. IACR ePrint 2020/1343.
\bibitem{Beu22} W.~Beullens. Breaking Rainbow takes a weekend on a laptop. \emph{CRYPTO} 2022. IACR ePrint 2022/214.
\bibitem{CGS14} P.~Campbell, M.~Groves, D.~Shepherd. Soliloquy: a cautionary tale. \emph{ETSI 2nd Quantum-Safe Crypto
  Workshop}, 2014.
\bibitem{CD23} W.~Castryck, T.~Decru. An efficient key recovery attack on SIDH. \emph{EUROCRYPT} 2023. IACR ePrint
  2022/975.
\bibitem{CIV16} W.~Castryck, I.~Iliashenko, F.~Vercauteren. Provably weak instances of Ring-LWE revisited.
  \emph{EUROCRYPT} 2016. IACR ePrint 2016/239.
\bibitem{CMP14} A.~Couvreur, I.~M\'arquez-Corbella, R.~Pellikaan. A polynomial time attack against algebraic geometry
  code based public key cryptosystems. \emph{IEEE ISIT} 2014. IACR ePrint 2014/064.
\bibitem{CDPR16} R.~Cramer, L.~Ducas, C.~Peikert, O.~Regev. Recovering short generators of principal ideals in
  cyclotomic rings. \emph{EUROCRYPT} 2016. IACR ePrint 2015/313.
\bibitem{DvW21} L.~Ducas, W.~van Woerden. NTRU fatigue: how stretched is overstretched? \emph{ASIACRYPT} 2021. IACR
  ePrint 2021/999.
\bibitem{ELOS15} Y.~Elias, K.~E.~Lauter, E.~Ozman, K.~E.~Stange. Provably weak instances of Ring-LWE. \emph{CRYPTO}
  2015. IACR ePrint 2015/106.
\bibitem{FGOPT13} J.-C.~Faug\`ere, V.~Gauthier-Uma\~na, A.~Otmani, L.~Perret, J.-P.~Tillich. A distinguisher for
  high-rate McEliece cryptosystems. \emph{IEEE Trans.\ Inf.\ Theory} 59(10) (2013). IACR ePrint 2010/331.
\bibitem{FOPT10} J.-C.~Faug\`ere, A.~Otmani, L.~Perret, J.-P.~Tillich. Algebraic cryptanalysis of McEliece variants
  with compact keys. \emph{EUROCRYPT} 2010, LNCS 6110, 279--298.
\bibitem{GifFLT} M.~Gifford. Fluid Lattice Topology: bridging PSL(2,7) vacuum geometry and non-associative post-quantum
  cryptography. Zenodo, 2026; the version with a notice withdrawing its security claims is
  doi:\doi{10.5281/zenodo.23351706}, on the record doi:\doi{10.5281/zenodo.20078312}.
\bibitem{GifNote} M.~Gifford. A rank collapse in a sedenion Module-LWE proposal, and why a uniform matrix does not
  repair it. Zenodo, 2026, version 2. doi:\doi{10.5281/zenodo.23352161} (all versions:
  doi:\doi{10.5281/zenodo.23270248}).
\bibitem{Kani97} E.~Kani. The number of curves of genus two with elliptic differentials. \emph{J.\ reine angew.\ Math.}
  485 (1997), 93--121.
\bibitem{KF17} P.~Kirchner, P.-A.~Fouque. Revisiting lattice attacks on overstretched NTRU parameters.
  \emph{EUROCRYPT} 2017.
\bibitem{LS15} A.~Langlois, D.~Stehl\'e. Worst-case to average-case reductions for module lattices.
  \emph{Des.\ Codes Cryptogr.} 75(3) (2015), 565--599.
\bibitem{MMPPW23} L.~Maino, C.~Martindale, L.~Panny, G.~Pope, B.~Wesolowski. A direct key recovery attack on SIDH.
  \emph{EUROCRYPT} 2023. IACR ePrint 2022/1026.
\bibitem{FIPS} NIST. FIPS 203: Module-lattice-based key-encapsulation mechanism standard; FIPS 204: Module-lattice-based
  digital signature standard. August 2024.
\bibitem{Pei16} C.~Peikert. How (not) to instantiate Ring-LWE. \emph{SCN} 2016. IACR ePrint 2016/351.
\bibitem{Regev} O.~Regev. On lattices, learning with errors, random linear codes, and cryptography.
  \emph{J.\ ACM} 56(6) (2009), article 34.
\bibitem{Rob23} D.~Robert. Breaking SIDH in polynomial time. \emph{EUROCRYPT} 2023. IACR ePrint 2022/1038.
\bibitem{SS92} V.~M.~Sidelnikov, S.~O.~Shestakov. On insecurity of cryptosystems based on generalized Reed--Solomon
  codes. \emph{Discrete Math.\ Appl.} 2(4) (1992), 439--444.
\end{thebibliography}

\end{document}
